
The classic function preamble, inserted by the compiler at the beginning of a function, can be seen in \cref{lst:classic_preamble}.
\begin{lstlisting}[caption={Classic function preamble supposing a function that uses 0x30 bytes of stack},label=lst:classic_preamble]
push rbp 
mov rbp, rsp
sub rsp, 0x30
\end{lstlisting}
Every function has its own \textbf{stack frame}, that is the stack zone where all the local informations used by the function 
lies. Specifically it contains local variables, parameters passed to the function, \dots
The stack frame is delimited by the value of two registers:
\begin{itemize}
    \item \textbf{Base Pointer} : point to the base of the activation frame (high addresses)
    \item \textbf{Stack Pointer} : point to the top of the stack (low addresses)
\end{itemize}
When a function is invoked, the stack frame of the new function must be allocated. This is done by the \cref{lst:classic_preamble}.
The effect on the stack could be seen in \cref{fig:preamble1}. If the function had a local variable, that variable 
would be referenced with respect to the base pointer (e.g. \texttt{[rbp - 0x8]}).
\begin{figure}[ht]
    \centering
    \includegraphics[width=\linewidth]{preamble1.png}
    \caption{Example of what happend during a function call}
    \label{fig:preamble1}
\end{figure}
    The base pointer in this case is called \textbf{frame pointer} and it is used as a stable reference to the stack frame.
Infact, during the execution of a function, the stack pointer could change 
(a \texttt{push} or a \texttt{pop} instruction occur), and this case the stack pointer would not be a valid reference to the 
variables anymore.\\
If the compiler knows a priori that the stack pointer will not change during function execution, could optimize the function 
preamble with a technique called \textbf{frame pointer omission}.
This technique consist in not use the base pointer as reference to the local variables, but using directly the stack 
pointer. This allow the function preamble to be shrink to the one shown in \cref{lst:fp_omission_preamble}
\begin{lstlisting}[caption={Function preamble in case of frame pointer omission supposing a function that uses 0x30 bytes of stack},label=lst:fp_omission_preamble]
sub rsp, 0x30
\end{lstlisting}
In this case the access to local variable is expressed with respect to the stack pointer 
(e.g. \texttt{[rsp + 0x8]}), as can be seen in \cref{fig:preamble2}.\\
\begin{figure}[ht]
    \centering
    \includegraphics[width=0.55\linewidth]{preamble2.png}
    \caption{Example of what happend during a function call in case of frame pointer omission}
    \label{fig:preamble2}
\end{figure}
\begin{center}
    \textbf{red zone}
\end{center}
The red zone is composed of 128 bytes above the rsp that are reserved by the ABI System V AMD64.
The user code that is compliant with the ABI could use this reserved memory to save temporary data without actually 
modify the stack pointer and allocate a stack frame.
Since the stack grows toward low addresses, this zone is overwritten by the stack frame of any function that is called from 
the actual one, so the red zone could be used especially for leaf functions (functions that does not call other functions).\\
The compiler could choose to bypass entirely the function premable and access the local variable with respect to 
the stack pointer (e.g.\texttt{[rsp -0x8]}) as shown in \cref{fig:preamble3}
\begin{figure}[ht]
    \centering
    \includegraphics[width=0.25\linewidth]{preamble3.png}
    \caption{Example of function call that does not allocate space of the stack bu access the variable saved on the red zone}
    \label{fig:preamble3}
\end{figure}